\chapter{下降技术与代数几何中的存在定理：Hilbert 概形}
\chapterintro{Techniques de construction et théorèmes d'existence en géométrie algébrique. IV : Les schémas de Hilbert}{Bourbaki，第221号报告；1961年5月}

\origpage{139}{221-01}
\section*{引言}
文献[2]的第一、二篇所介绍的技术，基本上不需要对所考察的\term{概形}{scheme}{schéma}作任何\term{射影性}{projectivity}{projectivité}假设。遗憾的是，目前这些技术还不能解决[2]第二篇所提出的存在性问题。本篇和下一篇将加入射影性假设，解决这些问题。所用方法具有典型的射影几何特征，实际上几乎不使用[2]第一、二篇的结果。本篇将构造\term{Hilbert 概形}{Hilbert scheme}{schéma de Hilbert}；正如[2]第二篇第2节所说，这一概形将取代\term{Chow 坐标}{Chow coordinates}{coordonnées de Chow}的使用。下一篇则把第三篇建立的概形取商理论与 Hilbert 概形理论结合起来；例如，我们将能够在相当一般的条件下，构造[2]第二篇第3节所定义的\term{Picard 概形}{Picard scheme}{schéma de Picard}。

可以说，现在我们已有一套大致令人满意的射影构造技术；尚缺的一环\footnote{见本篇末尾的补遗。}，是对射影群这类群的“无不动点”作用建立取商定理（参见[2]第三篇第8节）。在解析几何中，如果仅研究给定解析空间上的“射影”解析空间，情形甚至似乎还要好一些：对于解析空间，群以良好方式作用时，取商的困难就消失了。另一方面，无论代数几何还是解析几何，都还需要发展一种不依赖射影性假设的构造技术。

\section{层的有界族：稳定性}
设 $k$ 为域，$X$ 为 $k$ 上的\term{预概形}{prescheme}{préschéma}；为简便起见，假设它是\term{有限型}{of finite type}{de type fini}的。对每个扩张 $K/k$，由此得到 $K$ 上的预概形 $X_K=X\otimes_k K$。若 $\mathcal F$ 是 $X_K$ 上的\term{凝聚层}{coherent sheaf}{faisceau cohérent}，$K'$ 是 $K$ 的扩张，则 $\mathcal F\otimes_K K'=\mathcal F_{K'}$ 是 $X_K\otimes_K K'=X_{K'}$ 上的\term{拟凝聚层}{quasi-coherent sheaf}{faisceau quasi-cohérent}。现在，设 $K,K'$ 是 $k$ 的任意两个扩张，$\mathcal F$、$\mathcal F'$ 分别是 $X_K$、$X_{K'}$ 上的拟凝聚层。

\origpage{140}{221-02}
如果存在将 $K,K'$ 映入 $k$ 的同一个扩张 $K''$ 的 $k$-同态，使相应的 $\mathcal F_{K''}$ 与 $\mathcal F'_{K''}$ 在 $X_{K''}$ 上同构，就称 $\mathcal F$ 与 $\mathcal F'$ \term{等价}{equivalent}{équivalents}。这确实是一个等价关系；以下将研究关于此关系的层的等价类，以及这些等价类组成的集合。注意，若 $X$ 在 $k$ 上有限型，那么每个凝聚层的类，都可以由某个 $X_K$ 上的凝聚层定义，其中 $K$ 是 $k$ 的有限型扩张。因此，在定义凝聚层的类时，可以只考虑 $k$ 的\term{代数闭扩张}{algebraically closed extension}{extension algébriquement close}；还可以固定一个超越次数无限的代数闭扩张 $\Omega/k$。此时，$X_\Omega$ 上的两个凝聚层 $\mathcal F,\mathcal F'$ 等价，当且仅当存在 $\Omega$ 的一个 $k$-自同构 $\sigma$，使 $\mathcal F\otimes_k(\Omega,\sigma)$ 与 $\mathcal F'$ 同构。容易看出，这两种定义所得的凝聚层的类之间有自然的一一对应。

设 $E,E'$ 为 $X$ 上的两个凝聚层类集合。考虑所有形如 $\mathcal F\otimes\mathcal F'$ 的层的类，其中 $\mathcal F,\mathcal F'$ 定义在同一个 $X_K$ 上，且其类分别属于 $E,E'$。这样得到的凝聚层类集合记为 $E\otimes E'$。类似地定义 $\mathcal{T}or_i(E,E')$ 等。一般地，设运算 $\mathcal U$ 把同一个 $X_K$ 上的 $n$ 个凝聚层 $\mathcal F_1,\ldots,\mathcal F_n$，对应到 $X_K$ 上的一组凝聚层 $\mathcal U(\mathcal F_1,\ldots,\mathcal F_n)$，并具有与层同构以及\term{基变换}{base change}{changement de base}下的逆像相容这一显然应有的性质。那么，它也诱导一个仍记作 $\mathcal U$ 的运算，把 $n$ 个凝聚层类集合 $E_1,\ldots,E_n$ 对应到一个凝聚层类集合 $\mathcal U(E_1,\ldots,E_n)$。

本节的目的是定义一类称为\term{有界}{bounded}{limités}的层类集合，并证明最常用的运算 $\mathcal U$ 作用于有界集合后，所得集合仍然有界。

设 $X$ 是诺特预概形 $S$ 上的有限型预概形。对每个 $s\in S$，\term{纤维}{fibre}{fibre} $X_s$ 是剩余域 $k(s)$ 上的有限型预概形，因而可以按上述意义考虑 $X_s$ 上的凝聚层类。“$X/S$ 的某一纤维上的凝聚层类”及类似说法，由此都有了确定含义。同样，逐纤维地进行前述构造，可以考虑 $E\otimes E'$ 等运算：它们把若干个 $X/S$ 纤维上的凝聚层类集合，

\origpage{141}{221-03}
对应到另一个 $X/S$ 纤维上的凝聚层类集合。

\begin{statement}{定义 1.1}
设 $E$ 是 $X/S$ 的纤维上的凝聚层类的集合。若存在 $S$ 上的有限型预概形 $S'$，以及 $X'=X\times_S S'$ 上的凝聚层 $\mathcal F'$，使 $E$ 包含在 $\mathcal F'$ 所定义的 $X/S$ 纤维上的层类集合中，就称 $E$ 为\term{有界族}{bounded family}{famille limitée}。
\end{statement}
按照定义，后一个集合在每个 $s\in S$ 上，由层 $\mathcal F'\otimes_{S'}k(s')$ 的类构成，其中 $s'$ 遍历 $S'$ 中位于 $s$ 上方的点。这里 $k(s')$ 是 $k(s)$ 的扩张，而 $(X_s)_{k(s')}$ 与 $X'$ 在 $s'$ 处的纤维 $X'\otimes_{S'}k(s')=X'_{s'}$ 自然相同。因此也可以说，有界族就是包含在某个凝聚层的\term{代数族}{algebraic family}{famille algébrique}中的族，而该代数族由一个 $S$ 上的有限型预概形 $S'$ 参数化。

有限个有界族的并仍然有界：把控制这些族的代数族的参数预概形 $S_i$ 取不交并即可。基变换 $T\to S$ 把关于 $X/S$ 的有界族变成关于 $X_T/T$ 的有界族；若 $T\to S$ 满射，逆命题也成立。更一般地，只需它的像包含给定族 $E$ 在 $X/S$ 上实际涉及的所有点 $s$。这从理论上把有界族的判定归结到 $S$ 是整数环 $\mathbb Z$ 上有限型代数的谱的情形。

若 $E,E'$ 是关于 $X/S$ 的有界层类族，则 $E\otimes E'$ 也有界。事实上，若 $E,E'$ 分别被由 $T\to S$、$X_T$ 上的层 $\mathcal F$，以及 $T'\to S$、$X_{T'}$ 上的层 $\mathcal F'$ 定义的代数族控制，那么 $E\otimes E'$ 被如下代数族控制：取 $T''=T\times_ST'\to S$，并取 $X_{T''}$ 上的层 $\mathcal F''$ 为 $\mathcal F$ 与 $\mathcal F'$ 拉回的张量积。这个论证之所以成立，是因为关于 $\mathcal F,\mathcal F'$ 的张量积函子右正合，所以与基扩张相容，特别是与取纤维相容。对 $\mathcal{T}or_i(E,E')$、$\mathcal{H}om(E,E')$、$\mathcal{E}xt^i(E,E')$ 这类局部运算，上述论证不能直接照搬。不过，仍可仿照 $E\otimes E'$ 的情形，证明这些运算把有界集合变为有界集合；只需额外使用如下类型的结果（均包含在[3]第四章6.11中）：每个有界族 $E$ 都能由一个代数族控制，该代数族由 $X_T$ 上关于 $T$ \term{平坦}{flat}{plat}的凝聚层 $\mathcal F$ 定义，且 $T$ 在 $S$ 上有限型。为此要把原来的参数空间“分成若干块”。适当层的这些平坦性，

\origpage{142}{221-04}
保证了 $\mathcal{T}or_i(\mathcal F,\mathcal F')$ 等运算与任意基变换相容。同样的方法也适用于全局性质的运算，例如凝聚层沿\term{固有态射}{proper morphism}{morphisme propre}的直接像、导出直接像，以及相对于固有态射的全局 $\mathcal{E}xt$（参见[5]第三章第6节）等。这些运算都把有界层族变为有界层族。注意，在这些运算中，所取层所在的预概形也可以发生变化。

下面两个命题基本上也由同一种平坦性技术证明。关于有限型态射的纤维上的\term{准素分解}{primary decomposition}{décomposition primaire}，特别参见[5]第四章。

\begin{statement}{命题 1.2}
设 $E,E'$ 为 $X/S$ 纤维上的两个有界层类集合，并假设 $X$ 在 $S$ 上固有。则：
\begin{enumerate}[label=(\roman*)]
\item 所有同态 $\mathcal F\to\mathcal F'$ 的核、余核、像构成有界族，其中 $\mathcal F$ 的类属于 $E$，$\mathcal F'$ 的类属于 $E'$。
\item 由类分别属于 $E,E'$ 的层 $\mathcal F,\mathcal F'$ 构成的扩张层 $\mathcal F''$，组成有界族。
\end{enumerate}
\end{statement}
经过适当基变换，可以假设 $E,E'$ 分别由 $X_T/T$ 上的凝聚层 $\mathcal F,\mathcal F'$ 定义，且 $T$ 在 $S$ 上有限型。还可以假设满足某些平坦性条件，从而层
\[
f_{T*}\mathcal{H}om_{X_T}(\mathcal F,\mathcal F'),\qquad
\mathcal{E}xt^1_{f_T}(\mathcal F,\mathcal F')
\]
的形成与任意态射 $T'\to T$ 所给出的基变换相容。此外，也可假设这两个 $T$ 上的凝聚层都是\term{局部自由}{locally free}{localement libres}的。令 $T_0,T_1$ 为 $T$ 上的\term{向量丛}{vector bundle}{fibré vectoriel}，其截面芽层分别为上述两个层。于是自然地定义出 $X_{T_0}$ 上凝聚层的同态 $\mathcal F_{T_0}\to\mathcal F'_{T_0}$，以及 $X_{T_1}$ 上凝聚层的扩张
\[
0\longrightarrow\mathcal F_{T_1}\longrightarrow\mathcal F''
\longrightarrow\mathcal F'_{T_1}\longrightarrow0,
\]
它们都具有显然的泛性质。第二个层给出的代数族控制了(ii)所考虑的族。同样，前述同态的核、余核和像

\origpage{143}{221-05}
给出(i)所需的族；这里须假设余核关于 $T_0$ 平坦，而这一点仍可通过把 $T_0$ 分成若干块来实现。

\begin{statement}{命题 1.3}
设 $E$ 是 $X/S$ 纤维上的层类的有界族。那么，所有 $(\operatorname{Supp}\mathcal F)_{\mathrm{red}}$ 的结构层的类，也构成有界族；这里 $\mathcal F$ 是某个 $X_K$ 上的凝聚层，$K$ 代数闭，且 $\mathcal F$ 的类属于 $E$。
\end{statement}
这里 $(\operatorname{Supp}\mathcal F)_{\mathrm{red}}$ 表示赋予诱导\term{约化结构}{reduced structure}{structure réduite}的 $\mathcal F$ 的支集；换言之，其结构层是 $\mathcal O_{X_K}$ 对定义 $\operatorname{Supp}\mathcal F$ 的最大理想层的商。对于通过准素分解理论由 $\mathcal F$ 自然得到的其他层，也可以证明与(1.3)类似的结论。例如 $\mathcal F/\mathcal F_i$，其中 $\mathcal F_i$ 是 $\mathcal F$ 中对应于其支集各分支的准素子层，并在具有此性质的子层中最小；又如 $\mathcal O_{X_K}/\mathfrak p$，其中 $\mathfrak p$ 为 $\mathcal F$ 的\term{相伴素理想层}{associated prime ideal sheaf}{faisceau d'idéaux premier associé}；或 $\mathcal O_{X_K}/\mathfrak q$，其中 $\mathfrak q$ 是与 $\mathcal F$ 的支集某个分支相应的准素理想层。这里基域均取为代数闭域。

\section{有界族与 Hilbert 多项式}
以下假设 $X$ 在 $S$ 上射影，并给定一个\term{甚丰沛层}{very ample sheaf}{faisceau très ample} $\mathcal O_X(1)$。对任意 $s\in S$ 及剩余域 $k(s)$ 的扩张 $K$，在 $X_K$ 上考虑相应的层 $\mathcal O_{X_K}(1)$；它仍是甚丰沛的。

对 $X_K$ 上的每个凝聚层 $\mathcal F$，定义函数
\[
P_{\mathcal F}(n)=\chi(X_K,\mathcal F(n)),
\]
即 $\mathcal F(n)$ 在 $X_K$ 上的\term{Euler--Poincaré 特征}{Euler--Poincaré characteristic}{caractéristique d'Euler-Poincaré}。这是关于整数 $n$ 的一个多项式，称为 $\mathcal F$ 的\term{Hilbert 多项式}{Hilbert polynomial}{polynôme de Hilbert}。对充分大的 $n$，$P_{\mathcal F}(n)$ 就是 $H^0(X_K,\mathcal F(n))$ 的 $K$-维数，因为此时 $i>0$ 的上同调 $H^i(X_K,\mathcal F(n))$ 全部为零。

若 $\mathcal F$ 是 $X$ 上关于 $S$ 平坦的凝聚层，则在属于 $S$ 的同一连通分支的纤维 $X_s$ 上，诱导层 $\mathcal F_s$ 的 Hilbert 多项式彼此相同（[5]第三章第7节）。由此可知，对 $X$ 上任意凝聚层 $\mathcal F$，即使不作平坦性假设，所有 $\mathcal F_s$（$s\in S$）的 Hilbert 多项式也只有有限多个。

\origpage{144}{221-06}
另一方面，回顾一下：$X$ 上任意凝聚层 $\mathcal F$ 都同构于某个 $\mathcal O_X(-n)^N$ 的商层，只要 $n,N$ 充分大。因此，纤维上的诱导层 $\mathcal F_s$ 也都是相应纤维上 $\mathcal O(-n)^N$ 的商。

以上两点给出了下列定理的必要性。
\begin{statement}{定理 2.1}
设 $X$ 在诺特预概形 $S$ 上射影，$\mathcal O_X(1)$ 关于 $S$ 甚丰沛，$E$ 为 $X/S$ 纤维上的层类集合。$E$ 有界的充分必要条件是：\footnote{“充分必要”据作者\hyperref[orig:8]{C-06页勘误}修正；原报告此处漏印“充分”，下文已经明确说明还要证明充分性。}
\begin{enumerate}[label=(\alph*)]
\item 存在 $X$ 上的凝聚层 $\mathcal L$（可取为 $\mathcal O_X(-n)^N$ 的形式），使 $E$ 包含在形如 $\mathcal L_K$ 的层的凝聚商层的类族中。
\item 类属于 $E$ 的层 $\mathcal F$ 的 Hilbert 多项式 $P_{\mathcal F}$，都属于同一个有限的多项式集合。
\end{enumerate}
\end{statement}
尚须证明充分性；这将是一个更精细结果的特例。设 $\mathcal F$ 为域 $K$ 上有限型预概形上的凝聚模，$r$ 为整数。令 $N_r$ 为如下子模：它在一个开集上的截面，是 $\mathcal F$ 在该开集上支集维数小于 $r$ 的所有截面。因此，当 $r>\dim\operatorname{Supp}\mathcal F$ 时 $N_r=\mathcal F$，当 $r\leq0$ 时 $N_r=0$。由此得到 $\mathcal F$ 的一个有限递增滤过，其因子 $N_{r+1}/N_r$ 的相伴素闭链，恰好是 $\mathcal F$ 的维数为 $r$ 的相伴素闭链。\translatornote{原文把滤过因子印为 $N_r/N_{r+1}$。按刚才的定义有 $N_r\subset N_{r+1}$，故商的顺序应为 $N_{r+1}/N_r$；此处据定义改正。}记
\[
\mathcal F_{(r)}=\mathcal F/N_r.
\]
于是，$\mathcal F_{(r)}$ 的相伴素闭链恰好是 $\mathcal F$ 的维数不小于 $r$ 的相伴素闭链。特别地，$\mathcal F_{(r)}=\mathcal F$，当且仅当 $\mathcal F$ 的所有相伴素闭链的维数都不小于 $r$。

\begin{statement}{定理 2.2}
在(2.1)的假设下，设 $s$ 为整数，并假设 $E$ 满足条件(a)，以及条件(b)的以下弱化形式：

\noindent $(\mathrm b_s)$ 类属于 $E$ 的层 $\mathcal F$ 的 Poincaré 多项式 $P_{\mathcal F}$ 中，次数不小于 $s-1$ 的各项系数有界。

则当 $\mathcal F$ 的类遍历 $E$ 时，层 $\mathcal F_{(s)}$ 构成有界族。此外，$P_{\mathcal F}$ 中 $s-2$ 次项的系数有统一下界。
\end{statement}
\translatornote{原文 $(\mathrm b_s)$ 把次数条件印为 $\leq s-1$，这里改为 $\geq s-1$；后面的归纳证明和推论2.3均使用高次项系数有界。原印条件确实不足：在代数闭域上的 $X=\mathbb P^2$ 中，取 $s=1$、光滑 $d$ 次平面曲线 $C_d$（$d\geq3$），令 $g_d=(d-1)(d-2)/2$，再取长度 $g_d-1$ 的零维子概形 $Z_d$。层 $\mathcal F_d=\mathcal O_{C_d}\oplus\mathcal O_{Z_d}$ 是固定层 $\mathcal O_X^2$ 的商，且 $P_{\mathcal F_d}(n)=dn$ 的常数项恒为零；但 $\mathcal F_{d,(1)}=\mathcal O_{C_d}$ 的 Hilbert 多项式随 $d$ 取无限多值，故由(2.1)的必要性知其不构成有界族。这给出了原印不等号的反例。}

\origpage{145}{221-07}
由此得到：
\begin{statement}{推论 2.3}
假设类属于 $E$ 的层 $\mathcal F$ 的所有相伴素闭链的维数 $d$ 都满足 $s\leq d\leq r$。那么，(2.1)(b)只需对 $P_{\mathcal F}$ 中次数介于 $s-1$ 与 $r$ 之间的系数作要求。
\end{statement}
本节余下部分概述(2.2)的证明。关键是以下两个引理，其中第一个是众所周知的，概括了 Chow 坐标中有用的数学内容。

\begin{statement}{引理 2.4（Chow）}
考虑纤维 $X_K$ 的子概形 $Y$ 的结构层，其中 $K$ 是 $S$ 的某个剩余域的代数闭扩张，$Y$ 约化且所有不可约分支的维数同为 $r$，并把 $\mathcal O_Y$ 看作 $\mathcal O_{X_K}$ 的商层。若这些 $Y$ 的次数有界，则它们构成有界族。
\end{statement}
这里，$Y$ 的次数 $a$ 最方便的定义，是通过
\[
P_{\mathcal O_Y}(n)=a\frac{n^r}{r!}+\cdots
\]
的首项系数来给出。

\begin{statement}{引理 2.5}
设 $\mathcal L$ 是 $X$ 上的凝聚层，$E$ 是形如 $\mathcal L_K$ 的层的商层 $\mathcal F$ 的类的集合，其中 $K$ 为 $S$ 的剩余域的扩张。假设 $X/S$ 的纤维维数不超过 $r$，并记
\[
P_{\mathcal F}(n)=a_{\mathcal F}\frac{n^r}{r!}
+b_{\mathcal F}\frac{n^{r-1}}{(r-1)!}
+\text{次数小于 }r-1\text{ 的项}.
\]
则 $a_{\mathcal F}$ 有界，$b_{\mathcal F}$ 有统一下界。若 $b_{\mathcal F}$ 也有界，则 $\mathcal F_{(r)}$ 构成有界族。
\end{statement}
\begin{proof}
把 $S$ 换成覆盖它的一些子概形的不交并后，可以假设存在有限态射 $f:X\to\mathbb P^r_S$，使 $\mathcal O_X(1)$ 同构于 $\mathcal O_{\mathbb P^r_S}(1)$ 的逆像。因此，对 $X$ 上的每个凝聚层 $\mathcal F$，有 $P_{\mathcal F}=P_{f_*(\mathcal F)}$。另一方面，利用上一节的方法，容易验证：$X$ 上一族层 $\mathcal F$ 有界，当且仅当其直接像族 $f_*(\mathcal F)$ 有界。最后还有
\[
f_*(\mathcal F)_{(r)}=f_*(\mathcal F_{(r)}).
\]

\origpage{146}{221-08}
这就归结到 $X=\mathbb P^r_S$ 的情形。还可以假设，对适当的 $k,s$ 有 $\mathcal L=\mathcal O_{\mathbb P^r_S}(k)^s$。系数满足 $0\leq a_{\mathcal F}\leq s$，所以有界。由于 $a_{\mathcal F}$ 已经有界，$P_{\mathcal F}(n)$ 的 $n^{r-1}$ 项系数有统一下界（或有界），等价于
\[
P_{\mathcal F}(n-k)=P_{\mathcal F(-k)}(n)
\]
具有相同性质。于是还可以归结到 $\mathcal L=\mathcal O_{\mathbb P^r_S}^{s}$。

考虑正合列
\[
0\longrightarrow N_r\longrightarrow\mathcal F
\longrightarrow\mathcal F_{(r)}\longrightarrow0;
\]
从而 $P_{\mathcal F}=P_{\mathcal F_{(r)}}+P_{N_r}$。又因为 $P_{N_r}$ 的 $n^{r-1}$ 项系数非负（$\dim\operatorname{Supp}N_r\leq r-1$），故
\[
b_{\mathcal F_{(r)}}\leq b_{\mathcal F}.
\]
所以证明引理时，可以用 $\mathcal F_{(r)}$ 替代 $\mathcal F$，即假设所考虑的 $\mathcal L$ 的商层 $\mathcal F$ 都\term{无挠}{torsion-free}{sans torsion}。

由于 $\mathbb P^r_K$ 正规，$\mathcal F$ 在某个开集 $U=\mathbb P^r_K-Y$ 上局部自由，秩为 $a=a_{\mathcal F}$，且 $Y$ 的余维不小于2。因此 $\bigwedge^a\mathcal F$ 在 $U$ 上的限制是可逆层。又因 $\mathbb P^r_K$ 正则且 $Y$ 的余维不小于2，这个限制来自 $\mathbb P^r_K$ 上某个在同构意义下唯一的\term{可逆层}{invertible sheaf}{faisceau inversible}。后者必为 $\mathcal O_{\mathbb P^r_K}(d)$，其中整数 $d$ 唯一确定。由于 $\bigwedge^a\mathcal F$ 是
\[
\bigwedge^a\mathcal O_{\mathbb P^r_K}^{s}\simeq
\mathcal O_{\mathbb P^r_K}^{N},\qquad N=\binom{s}{a},
\]
的商，它有 $N$ 个自然截面，因此给出

\origpage{147}{221-09}
$\mathcal O_{\mathbb P^r_K}(d)$ 在 $U$ 上的截面。这些截面又是该可逆层在 $\mathbb P^r_K$ 上的截面 $s_i$（$1\leq i\leq N$）的限制，因为 $\mathbb P^r_K$ 正规且 $Y$ 的余维不小于2。它们在 $U$ 上生成 $\mathcal O_{\mathbb P^r_K}(d)$，所以不全为零，从而 $d\geq0$。另一方面，容易计算出
\[
b_{\mathcal F}=a_{\mathcal F}\frac{r+1}{2}+d.
\]
特别地，$b_{\mathcal F}\geq0$，因而有统一下界。它有界当且仅当 $d$ 有界。下面说明，此时 $\mathcal F$ 构成有界族。

可以固定 $a_{\mathcal F}=a$、$b_{\mathcal F}=b$，于是 $d$ 也固定。给定这 $N$ 个截面 $s_i$，即给定同态
\[
s:\bigwedge^a\mathcal L_K\longrightarrow\mathcal O_{\mathbb P^r_K}(d),
\]
就可以把 $\mathcal F$ 恢复为下列复合同态的余像：
\[
\begin{aligned}
\mathcal L_K&\longrightarrow
\mathcal{H}om\!\left(\bigwedge^{a-1}\mathcal L_K,
\bigwedge^a\mathcal L_K\right)\\
&\longrightarrow\mathcal{H}om\!\left(\bigwedge^{a-1}\mathcal L_K,
\mathcal O_{\mathbb P^r_K}(d)\right).
\end{aligned}
\]
第一个箭头是外积给出的自然同态，第二个由 $s$ 导出。由(1.2)(i)即得结论。
\end{proof}

结合上述两个引理，可证明：
\begin{statement}{引理 2.6}
在(2.1)的基本假设下，设每个 $\mathcal F$ 均有
\[
P_{\mathcal F}(n)=a_{\mathcal F}\frac{n^r}{r!}
+b_{\mathcal F}\frac{n^{r-1}}{(r-1)!}
+\text{次数小于 }r-1\text{ 的项},
\]
且系数 $a_{\mathcal F}$ 有界。则 $b_{\mathcal F}$ 有统一下界。若 $b_{\mathcal F}$ 有界，则 $\mathcal F_{(r)}$ 构成有界族。
\end{statement}
可以假设各层 $\mathcal F$ 的基域 $K$ 均代数闭。给每个 $\operatorname{Supp}\mathcal F_{(r)}$（即 $\operatorname{Supp}\mathcal F$ 的所有 $r$ 维分支之并）赋予诱导的约化结构。它们的次数被 $a_{\mathcal F}$ 控制，所以根据(2.4)，$\operatorname{Supp}\mathcal F_{(r)}$ 构成有界族。此外，$\mathcal F_{(r)}$ 在其支集的每个分支处的长度不超过 $a_{\mathcal F}$。若 $\mathcal J_{\mathcal F}$ 是定义 $\operatorname{Supp}\mathcal F_{(r)}$ 的理想，则 $\mathcal F_{(r)}$

\origpage{148}{221-10}
可以看作由 $\mathcal J_{\mathcal F}^{a_{\mathcal F}}$ 定义的 $X$ 的子概形 $Y_{\mathcal F}$ 上的模。如前一引理一样，还可以归结到 $\mathcal F=\mathcal F_{(r)}$ 的情形，于是 $\mathcal F$ 来自 $Y_{\mathcal F}$ 上的模。这些 $Y_{\mathcal F}$ 对应于 $\mathcal O_{X_K}$ 的商模的有界族，因而来自某个 $X\times_ST$ 的闭子概形 $Y$。现在对 $Y/T$ 与 $\mathcal L\otimes_{\mathcal O_X}\mathcal O_Y$ 应用(2.5)，便得结论。

现在对 $\dim\operatorname{Supp}\mathcal F$ 的上界 $r$ 归纳，证明(2.2)。\translatornote{归纳的边界情形须先分出：若 $r<s$，则 $\mathcal F_{(s)}=0$，故其有界性显然。若 $r<s-2$，所求 $s-2$ 次项系数为零；若 $r=s-2$，它是非负的最高项系数；若 $r=s-1$，条件$(\mathrm b_s)$控制最高项，再由引理2.6得到下一项的统一下界。因此，下面使用 $\mathcal F_{(r)}$ 有界的归纳步骤只需在 $r\geq s$ 时进行；此时$(\mathrm b_s)$确实控制了应用2.6所需的最高两项。}当 $r<0$ 时结论显然；设 $r\geq0$，且对所有 $r'<r$ 已证明结论。由(2.6)，$\mathcal F_{(r)}$ 构成有界族；再由(1.2)(i)，同态 $\mathcal L_K\to\mathcal F_{(r)}$ 的核也构成有界族。因此存在 $X$ 上的凝聚模 $\mathcal L'$，使这些核，以及
\[
N_r(\mathcal F)=\ker(\mathcal F\to\mathcal F_{(r)}),
\]
都是 $\mathcal L'_K$ 的商模。由于 $\mathcal F_{(r)}$ 有界，多项式 $P_{\mathcal F_{(r)}}$ 有界，而等式
\[
P_{\mathcal F}=P_{\mathcal F_{(r)}}+P_{N_r}
\]
表明 $P_{N_r}$ 满足与 $P_{\mathcal F}$ 相同的条件 $(\mathrm b_s)$。所以 $N_r$ 满足(a)与$(\mathrm b_s)$；由归纳假设，$(N_r)_{(s)}$ 有界。又 $\mathcal F_{(s)}$ 是 $\mathcal F_{(r)}$ 被 $(N_r)_{(s)}$ 的扩张，因此由(1.2)(ii)，$\mathcal F_{(s)}$ 有界。至于(2.2)的最后一个断言，先由(1.2)(i)知 $\mathcal L_K\to\mathcal F_{(s)}$ 的核构成有界族，因此它们的商 $N_s=\ker(\mathcal F\to\mathcal F_{(s)})$ 都是某个固定凝聚层的基扩张的商。又由条件$(\mathrm b_s)$及 $\mathcal F_{(s)}$ 有界，$P_{N_s}=P_{\mathcal F}-P_{\mathcal F_{(s)}}$ 的 $n^{s-1}$ 项系数有界。对 $N_s$ 应用引理2.6，便知其下一项系数有统一下界，从而 $P_{\mathcal F}$ 的 $s-2$ 次项系数也有统一下界。证明完成。\translatornote{原文此处直接断言 $N_s$ 有界，这一般不成立；例如 $X=\mathbb P^2$、$s=2$，取 $\mathcal F=\mathcal O_X\oplus\mathcal O_Z$，其中 $Z$ 是任意长度的零维子概形。这些层都是固定 $\mathcal O_X^2$ 的商，高次系数满足条件$(\mathrm b_s)$，且 $\mathcal F_{(2)}=\mathcal O_X$，但 $N_2=\mathcal O_Z$ 并不有界。正文将这一处简写补为实际所需的较弱结论：$N_s$ 有固定商源，且其最高可能次数 $s-1$ 的系数有界；这已经足以应用引理2.6，原定理不受影响。}

\section{Hilbert 概形：定义与存在定理}
设 $X$ 是预概形 $S$ 上的预概形，$\mathcal F$ 是 $X$ 上的拟凝聚模。记
\[
\operatorname{Quot}(\mathcal F/X/S)
\]
为 $\mathcal F$ 的所有关于 $S$ 平坦的拟凝聚商模组成的集合。给定基变换态射 $S'\to S$，令 $X'=X\times_SS'$、$\mathcal F'=\mathcal F\otimes_{\mathcal O_S}\mathcal O_{S'}$。于是 $X'$ 是 $S'$ 上的预概形，并带有拟凝聚模

\origpage{149}{221-11}
$\mathcal F'$，因而可考虑 $\operatorname{Quot}(\mathcal F'/X'/S')$。定义
\[
\underline{\operatorname{Quot}}_{\mathcal F/X/S}(S')
=\operatorname{Quot}(\mathcal F'/X'/S'),\qquad X'=X\times_SS'.
\]
若 $S''\to S'$ 是 $S$-态射，则 $X''=X\times_SS''$ 同构于 $X'\times_{S'}S''$，而 $\mathcal F''$ 同构于 $\mathcal F'\otimes_{\mathcal O_{S'}}\mathcal O_{S''}$。逆像函子
\[
\mathcal G'\longmapsto\mathcal G'\otimes_{\mathcal O_{S'}}\mathcal O_{S''}
\]
从 $X'$ 上拟凝聚模范畴到 $X''$ 上拟凝聚模范畴，是右正合的，并把 $S'$-平坦模变为 $S''$-平坦模。因此得到自然映射
\[
\operatorname{Quot}(\mathcal F'/X'/S')\longrightarrow
\operatorname{Quot}(\mathcal F''/X''/S'').
\]
所以 $\underline{\operatorname{Quot}}_{\mathcal F/X/S}$ 是关于 $S$ 上预概形 $S'$ 的集合值\term{反变函子}{contravariant functor}{foncteur contravariant}。以下假设 $X$ 在诺特预概形 $S$ 上射影，$\mathcal F$ 凝聚；为简便起见，只考虑局部诺特的 $S$ 上预概形 $S'$。

\begin{statement}{定理 3.1}
在上述假设下，局部诺特 $S$-预概形范畴上的反变函子 $\underline{\operatorname{Quot}}_{\mathcal F/X/S}$，可由一个 $S$-预概形 $\operatorname{Quot}_{\mathcal F/X/S}$ 表示。后者是一列射影 $S$-概形的不交并；特别地，它在 $S$ 上局部有限型。
\end{statement}
我们将如下得到这种分解。设 $\mathcal O_X(1)$ 是 $X$ 上关于 $S$ 甚丰沛的可逆层。对每个有理系数多项式 $P(n)$，令 $\operatorname{Quot}^{P}(\mathcal F/X/S)$ 为 $\operatorname{Quot}(\mathcal F/X/S)$ 中的如下部分：由 $\mathcal F$ 的凝聚商层 $\mathcal G$ 构成，$\mathcal G$ 在 $S$ 上平坦，且在每个 $s\in S$ 处的 Hilbert 多项式都等于 $P$。再令
\[
\underline{\operatorname{Quot}}^{P}_{\mathcal F/X/S}(S')
=\operatorname{Quot}^{P}(\mathcal F'/X'/S').
\]
这样得到 $\underline{\operatorname{Quot}}_{\mathcal F/X/S}$ 的一个\term{子函子}{subfunctor}{sous-foncteur}。第2节回顾的 Hilbert 多项式不变性说明：$\underline{\operatorname{Quot}}_{\mathcal F/X/S}$ 可表，当且仅当所有 $\underline{\operatorname{Quot}}^{P}_{\mathcal F/X/S}$ 都可表；此时表示前者的 $S$-预概形 $\operatorname{Quot}_{\mathcal F/X/S}$，同构于表示这些子函子的 $\operatorname{Quot}^{P}_{\mathcal F/X/S}$ 的不交并。因此，(3.1)将由下列定理推出。

\origpage{150}{221-12}
\begin{statement}{定理 3.2}
沿用上述记号，函子 $\underline{\operatorname{Quot}}^{P}_{\mathcal F/X/S}$ 可由一个射影 $S$-预概形 $\operatorname{Quot}^{P}_{\mathcal F/X/S}$ 表示。
\end{statement}
本节以下证明(3.2)。设 $\nu$ 为整数。对每个 $S$ 上的 $S'$，令 $A_\nu(S')$ 为满足下列条件的凝聚商层 $\mathcal G=\mathcal F'/\mathcal H$ 所组成的集合；其中 $\mathcal F'=\mathcal F\otimes_{\mathcal O_S}\mathcal O_{S'}$，$\mathcal G$ 在 $S'$ 上平坦，且属于下述 $A(S')=\underline{\operatorname{Quot}}^P_{\mathcal F/X/S}(S')$：
\begin{align*}
\text{a.}\quad &R^if'_*(\mathcal G(n))=0 &&(i>0,\ n\geq\nu),\\
\text{b.}\quad &R^if'_*(\mathcal H(n))=0 &&(i>0,\ n\geq\nu),\\
\text{c.}\quad &f'_*(\mathcal H(\nu+k))=\mathcal S'_k\,f'_*(\mathcal H(\nu)) &&(k\geq0).
\end{align*}
最后一个关系中，假定 $X$ 写成 $S$ 上的拟凝聚分次代数 $\mathcal S_*$ 的齐次素谱；$\mathcal S_*$ 非负分次，并由 $\mathcal S_1$ 生成。记 $\mathcal S'_* =\mathcal S_*\otimes_{\mathcal O_S}\mathcal O_{S'}$，于是 $X'$ 是 $\mathcal S'_*$ 的齐次素谱。为证明(3.2)，也容易归结到 $X=\mathbb P^r_S$、$\mathcal F=\mathcal O_{\mathbb P^r_S}^{N}$ 的情形：$S$ 可以由开集 $U$ 覆盖，使 $X|_U$ 是某个 $\mathbb P^r_U$ 的闭子概形，且 $\mathcal O_X(1)$ 由 $\mathcal O_{\mathbb P^r_U}(1)$ 诱导。在这一情形，$\mathcal F$ 在 $S$ 上平坦，所以前面的 $\mathcal H$ 也在 $S'$ 上平坦。

由 Künneth 关系（[5]第三章第7节）以及(b)，条件(a)、(b)在基变换下稳定，并且蕴含：当 $n\geq\nu$ 时，$f'_*(\mathcal G(n))$ 与 $f'_*(\mathcal H(n))$ 的形成都与基扩张相容。因此(c)也在基扩张下稳定。换言之，$A_\nu$ 是关于 $S'$ 的反变函子，更准确地说，是 $A$ 的子函子。让 $\nu$ 变化，就得到 $A(S')$ 中的一列递增子集 $A_\nu(S')$；由 Serre 的一个著名定理（[5]第三章第2节），它们的并为 $A(S')$。

还要注意：若 $\mathcal G$ 是 $\mathcal F$ 的 $S$-平坦凝聚商层，且对某个 $s\in S$，基变换 $\operatorname{Spec}k(s)\to S$ 所得商层 $\mathcal G_s$ 满足(a)、(b)、(c)，即属于 $A_\nu(\operatorname{Spec}k(s))$，则存在 $s$ 的开邻域 $U$，使这些条件对 $\mathcal G|_{f^{-1}(U)}$ 仍成立，即该商层属于 $A_\nu(U)$。对(a)、(b)，这是“全纯函数定理”（[5]第三章第7节）的结果；对(c)，则利用 Nakayama 引理以及如下事实：只要 $n$ 充分大，对一切 $k\geq0$ 总有
\[
f_*(\mathcal H(n+k))=\mathcal S_k f_*(\mathcal H(n))
\]

\origpage{151}{221-13}
（参见[5]第三章第2节）。由这些观察可知（比较[4]第四篇）：函子 $A$ 可表的充分必要条件，是所有 $A_\nu$ 都可表；此时表示 $A$ 的 $S$-预概形 $Q$ 是一列递增开集 $Q_\nu$ 的并，且 $Q_\nu$ 表示 $A_\nu$。

记
\[
M_* =\bigoplus_{n\geq0}f_*(\mathcal F(n))=\mathcal S_*^N.
\]
于是
\[
M'_*=M_*\otimes_{\mathcal O_S}\mathcal O_{S'}
=\bigoplus_{n\geq0}f'_*(\mathcal F'(n))=(\mathcal S'_*)^N.
\]
由(a)以及[5]第三章第7节，得到：

\noindent(a$'$) 当 $n\geq\nu$ 时，$f'_*(\mathcal G(n))$ 局部自由，秩为 $P(n)$。

\noindent由(b)的 $i=1$ 情形，还得到：

\noindent(a$''$) $f'_*(\mathcal G(n))$ 是 $M'_n$ 的商模。

给定 $n=\nu$ 时的这个商模，由(c)可确定所有 $n\geq\nu$ 时 $M'_n$ 的子模 $f'_*(\mathcal H(n))$，因而确定 $\mathcal H$，从而确定 $\mathcal G$。于是得到单射
\[
A_\nu(S')\longrightarrow\operatorname{Grass}_{P(\nu)}(M'_\nu),
\]
其目标是 $M'_\nu$ 的局部自由秩 $P(\nu)$ 商模的集合。这给出函子间的同态
\[
i_\nu:A_\nu(S')\longrightarrow
\underline{\operatorname{Grass}}_{P(\nu)}(M_\nu)(S').
\]
右侧函子由\term{Grassmann 概形}{Grassmann scheme}{schéma grassmannien} $\operatorname{Grass}_{P(\nu)}(M_\nu)$ 表示（比较[4]第五篇），而后者在 $S$ 上射影。我们断言：

\begin{statement}{引理 3.3}
函子 $A_\nu$ 可表，且表示同态 $i_\nu$ 的态射 $Q_\nu\to\operatorname{Grass}_{P(\nu)}(M_\nu)$ 是\term{浸入}{immersion}{immersion}。特别地，$Q_\nu$ 在 $S$ 上\term{拟射影}{quasi-projective}{quasi-projectif}。
\end{statement}
这个断言等价于如下说法（比较[4]第四篇）：给定 $M'_\nu$ 的一个局部自由秩 $P(\nu)$ 商模 $N$，则存在 $S'$ 的一个子预概形 $Z$，使对每个局部诺特的 $S'$-预概形 $T'$，

\origpage{152}{221-14}
$N$ 在 $T'$ 上的逆像属于 $i_\nu A_\nu(T')$，当且仅当 $T'\to S'$ 经子预概形 $Z$ 分解。更换记号后可以假设 $S'=S$，即已给定由子模 $R_\nu$ 定义的 $M_\nu$ 的商模 $N_\nu$。它来自 $A(S)$ 的一个元素的充分必要条件，是满足下列两点：
\begin{enumerate}[label=(\roman*)]
\item 对每个 $k\geq0$，$M_{\nu+k}/\mathcal S_kR_\nu$ 局部自由，秩为 $P(\nu+k)$。
\item $M_*$ 的分次子模 $R_*=\bigoplus_{k\geq0}\mathcal S_kR_\nu$ 定义出 $\mathcal F$ 的子层 $\mathcal H$（[5]第二章第9节）；$\mathcal H$ 与商层 $\mathcal G=\mathcal F/\mathcal H$ 满足前面的条件(a)、(b)。
\end{enumerate}
这些条件显然必要。反过来，若它们成立，则(ii)中定义的 $\mathcal G$ 同构于分次 $\mathcal S_*$-模
\[
N_*=\bigoplus_{k\geq0}M_{\nu+k}/\mathcal S_kR_\nu
\]
的关联层，因而在 $S$ 上平坦：它的各茎是 $N_*$ 对 $\mathcal S_*$ 的齐次素理想局部化后所得模的直和因子。由(i)，它的 Hilbert 多项式为 $P$。结合(ii)，可知(a$'$)、(a$''$)成立。因此，当 $n\geq\nu$ 时，自然同态 $N_n\to f_*(\mathcal G(n))$ 是同秩局部自由模之间的满射，故为同构。所以对每个 $k\geq0$ 都有
\[
f_*(\mathcal H(\nu+k))=\mathcal S_kR_\nu.
\]
这证明 $\mathcal G\in A_\nu(S)$，且 $N_\nu=M_\nu/R_\nu$ 正是 $\mathcal G$ 所定义的 Grassmann 点。

这个判别准则(i)、(ii)在基变换 $S'\to S$ 后仍适用。我们先证明：基变换后条件(i)成立，等价于 $S'\to S$ 经 $S$ 的某个子预概形 $Z$ 分解。一旦得到这个结果，以 $Z$ 替代 $S$，即可假设条件(i)在 $S$ 上已经成立。由于它在基变换下稳定，剩下的只是表达条件(ii)。令 $U$ 为满足下列性质的点 $s\in S$ 的集合：对所有 $n\geq\nu$，$\mathcal G(n)$ 与 $\mathcal H(n)$ 在纤维 $X_s$ 上诱导的层的正次上同调均为零。前面已经指出 $U$ 是开集，且基变换后条件(ii)成立，当且仅当 $S'\to S$ 经 $U$ 分解。这就证明了引理3.3。因此，只需再证明：

\begin{statement}{引理 3.4}
设 $S$ 为局部诺特预概形，带有一个由 $\mathcal S_1$ 生成的非负分次拟凝聚代数 $\mathcal S_*$；$M_*$ 为有限型分次拟凝聚 $\mathcal S_*$-模，$P$ 为有理系数多项式，$\nu$ 为整数。则存在 $S$ 的一个子预概形 $Z$（显然唯一），

\origpage{153}{221-15}
具有如下性质：对每个 $S$ 上的预概形 $S'$，所有 $n\geq\nu$ 的 $M_n\otimes_{\mathcal O_S}\mathcal O_{S'}$ 都局部自由且秩为 $P(n)$，当且仅当 $S'\to S$ 经 $Z$ 分解。
\end{statement}
显然可以假设 $S$ 仿射，从而诺特。此时有：

\begin{statement}{引理 3.5}
对每个整数 $N\geq\nu$，令 $U_N$ 为由满足
\[
\dim_{k(s)}(M_n\otimes_{\mathcal O_S}k(s))\leq P(n)
\quad(\nu\leq n\leq N)
\]
的点 $s\in S$ 组成的开集。则递减开集列 $U_N$ 最终稳定。
\end{statement}
由[3]第四章，$S$ 可分割为有限多个约化子预概形 $S_i$，使每个 $M_*\otimes_{\mathcal O_S}\mathcal O_{S_i}$ 在 $S_i$ 上平坦。因此可假设 $M_*$ 平坦，从而各 $M_n$ 平坦。最后，还可假设 $S$ 连通。于是存在整数 $n_0$ 及多项式 $Q$，使
\[
\dim_{k(s)}(M_n\otimes_{\mathcal O_S}k(s))=Q(n)
\quad(n\geq n_0)
\]
（[5]第三章第7节）。先设 $P<Q$，即充分大 $n$ 时 $P(n)<Q(n)$。则 $N$ 充分大时显然 $U_N=\varnothing$，所以 $U_N$ 稳定。在其余情形，充分大 $n$ 时 $P(n)\geq Q(n)$，因而 $N\geq n_0$ 后有 $U_N=U_{n_0}$，开集列也稳定。\footnote{此处分情况以及下面公式中的非严格不等号，采用作者\hyperref[orig:8]{C-06页勘误}。}

特别地，所有满足
\[
\dim_{k(s)}(M_n\otimes_{\mathcal O_S}k(s))\leq P(n)
\quad\text{对所有 }n\geq\nu\tag{*}
\]
的点组成的集合 $U_\infty$，是 $U_N$ 的交，因而是开集。为证明(3.4)，可以用开集 $U_\infty$ 替代 $S$，这样就归结到每个 $s\in S$ 都满足(*)的情形。此外：

\begin{statement}{引理 3.6}
设 $M$ 是局部诺特预概形 $S$ 上的凝聚模，$r$ 为整数。则存在 $S$ 的一个子预概形 $Z$（显然唯一），具有如下性质：对每个 $S$ 上的 $S'$，$M\otimes_{\mathcal O_S}\mathcal O_{S'}$ 局部自由且秩为 $r$，当且仅当 $S'\to S$ 经 $Z$ 分解。若对每个 $s$ 都有 $\dim_{k(s)}(M\otimes_{\mathcal O_S}k(s))\leq r$，则 $Z$ 是 $S$ 的闭子预概形。
\end{statement}

\origpage{154}{221-16}
事实上，仿照前面的论证，可把 $S$ 换成不等式成立的点组成的开集，从而假设该不等式在每个 $s\in S$ 上成立。局部地，可以令 $M$ 出现在一个正合列
\[
\mathcal O_S^q\longrightarrow\mathcal O_S^r\longrightarrow M\longrightarrow0
\]
中。对 $S'$ 的要求就等价于：相应正合列
\[
\mathcal O_{S'}^q\longrightarrow\mathcal O_{S'}^r\longrightarrow M'\longrightarrow0
\]
中第二个箭头为同构，也就是第一个箭头为零。因此，以同态 $\mathcal O_S^q\to\mathcal O_S^r$ 的矩阵各系数所生成的理想定义闭子预概形 $Z$，它就具有所需性质。

现在回到(3.4)的证明。由(3.6)，对 $M_n$ 与整数 $r=P(n)$，得到 $S$ 的闭子概形 $Z_n$；再令 $Z'_N$ 为 $\nu\leq n\leq N$ 时所有 $Z_n$ 的概形论交。$Z'_N$ 构成一列递减的闭子概形，故最终稳定。取 $N$ 充分大时的稳定值为 $Z$，它就是(3.4)所要求的子预概形。这就完成了(3.4)，从而也完成了(3.3)的证明。

至此已经证明，$\underline{\operatorname{Quot}}^P_{\mathcal F/X/S}$ 可由一个 $S$-预概形 $Q$ 表示，且 $Q$ 是一列在 $S$ 上拟射影的开子预概形 $Q_\nu$ 的递增并。要进一步推进，需要调用定理2.1。由该定理容易推出 $Q$ \term{拟紧}{quasi-compact}{quasi-compact}：它是某个 $S$ 上有限型的参数预概形 $S'$ 的像，$S'$ 参数化具有 Hilbert 多项式 $P$ 的 $\mathcal F_K$ 的商层。因此 $Q$ 等于其中一个 $Q_\nu$，从而在 $S$ 上拟射影。要证明射影性，只需再证明 $Q$ 在 $S$ 上固有；利用[5]第二章7.3.8形式的\term{固有性的赋值判别法}{valuative criterion of properness}{critère valuatif de propreté}，只需验证：

\begin{statement}{引理 3.7}
设 $S$ 是\term{离散赋值环}{discrete valuation ring}{anneau de valuation discrète}的谱，$s$ 为其泛点；$X$ 为 $S$ 上的预概形，$\mathcal F$ 为 $X$ 上的拟凝聚模，$\mathcal G_s$ 为 $X_s$ 上的 $\mathcal F_s=\mathcal F\otimes_{\mathcal O_S}k(s)$ 的拟凝聚商模。则存在唯一的 $\mathcal F$ 的拟凝聚商模 $\mathcal G$，它在 $S$ 上平坦，且在 $X_s$ 上的限制为 $\mathcal G_s$。
\end{statement}

\origpage{155}{221-17}
事实上，若 $\mathcal G_s=\mathcal F_s/\mathcal H_s$，只需取 $\mathcal F$ 中诱导出 $\mathcal H_s$ 的最大子层 $\mathcal H$（[5]第一章9.4.2），然后令 $\mathcal G=\mathcal F/\mathcal H$。容易验证这个层满足要求。

定理3.2，以及由它推出的(3.1)，至此全部得证。证明同时还给出：

\begin{statement}{命题 3.8}
在(3.2)的条件下，令 $Q=\operatorname{Quot}^P_{\mathcal F/X/S}$、$X_Q=X\times_SQ$、$\mathcal F_Q=\mathcal F\otimes_{\mathcal O_S}\mathcal O_Q$。设 $\mathcal G$ 为 $\mathcal F_Q$ 的凝聚商层，在 $Q$ 上平坦，相对 Hilbert 多项式为 $P$，且 $(Q,\mathcal G)$ 表示函子 $\underline{\operatorname{Quot}}^P_{\mathcal F/X/S}$。则存在整数 $\nu$，使所有 $n\geq\nu$ 的 $(f_Q)_*(\mathcal G(n))$ 都是 $Q$ 上局部自由秩 $P(n)$ 的模，并且关于 $S$ 甚丰沛；也就是说，它定义了 $Q$ 到 $S$ 上某个 Grassmann 概形 $\operatorname{Grass}_{P(n)}(M)$ 的浸入。特别地，$n\geq\nu$ 时，$Q$ 上的层
\[
\bigwedge^{P(n)}(f_Q)_*(\mathcal G(n))
\]
是关于 $S$ 甚丰沛的可逆层。
\end{statement}
事实上，可以像证明(3.2)时一样，归结到 $\mathcal F$ 在 $S$ 上平坦的情形；然后只需取一个使 $A_\nu=A$ 的整数 $\nu$ 即可，这里使用前面的记号。

(3.2)最重要的应用是 $\mathcal F=\mathcal O_X$ 的情形。此时记
\[
\operatorname{Quot}_{\mathcal O_X/X/S}=\operatorname{Hilb}_{X/S},\qquad
\operatorname{Quot}^P_{\mathcal O_X/X/S}=\operatorname{Hilb}^P_{X/S},
\]
因而有分解
\[
\operatorname{Hilb}_{X/S}=\coprod_P\operatorname{Hilb}^P_{X/S}.
\]
按定义，$\operatorname{Hilb}_{X/S}$ 表示函子 $\underline{\operatorname{Hilb}}_{X/S}$，其在 $S'$ 上的值是 $X'=X\times_SS'$ 中所有在 $S'$ 上平坦的闭子预概形的集合；$\operatorname{Hilb}^P_{X/S}$ 则表示由 Hilbert 多项式为给定 $P$ 的闭子预概形构成的子函子。它们分别称为 $X$ 在 $S$ 上的 Hilbert 预概形，以及指标为 $P$ 的 Hilbert 预概形。这个名称反映了 Hilbert 多项式在理论中的作用。它们与经典 Chow 簇（用来参数化闭链，而非簇）在性质上的区别，正如一个簇的闭链类的 Chow 环，

\origpage{156}{221-18}
与该簇的层类环之间的区别；后一个环就是 Riemann--Roch 定理[1]中出现的环。事实上，当 $X=\mathbb P^r_S$ 且 $S$ 为域的谱时，给定 $X$ 上凝聚模 $\mathcal F$ 的 Hilbert 多项式，也就等价于给定 $\mathcal F$ 的 Chern 类，或者等价于给定 $\mathcal F$ 在 $X$ 的凝聚层类环中的类。\translatornote{这里“给定 Chern 类”须连同秩一起理解，或改用包括零次部分的 Chern 特征。若只给出通常的正次 Chern 类，$\mathcal O_X$ 与 $\mathcal O_X^2$ 的总 Chern 类都等于1，但 Hilbert 多项式分别为 $\binom{n+r}{r}$ 与 $2\binom{n+r}{r}$，并不能互相确定。把秩补入后，射影空间的 Riemann--Roch 公式及其 $K$ 群的无挠性给出文中所述对应。}

\subsection*{注 3.9}
还应注意，$\operatorname{Quot}_{\mathcal F/X/S}$ 与 $\operatorname{Quot}^P_{\mathcal F/X/S}$ 的构造，已经归结到 $X=\mathbb P^r_S$、$\mathcal F=\mathcal O_X^N$、$\mathcal O_X(1)$ 为通常的甚丰沛层的情形。确切地说，一般的 $\operatorname{Quot}^P_{\mathcal F/X/S}$ 都可实现为这些特殊情形中的闭子预概形。又因 $\operatorname{Quot}^P_{\mathcal F/X/S}$ 的形成显然与基变换 $S'\to S$ 相容，还可进一步归结到 $S=\operatorname{Spec}\mathbb Z$。最终，我们只需研究 $\mathbb Z$ 上的射影概形
\[
\mathfrak Q^P_{r,N}
=\operatorname{Quot}^P_{\mathcal O_{\mathbb P^r_{\mathbb Z}}^N/
\mathbb P^r_{\mathbb Z}/\operatorname{Spec}\mathbb Z},
\]
特别是绝对 Hilbert 概形
\[
\operatorname{Hilb}^P_r=\mathfrak Q^P_{r,1}.
\]
详细研究这些概形会很有意思，首先就是确定其连通分支：它们是否连通？其次是确定其不可约分支；由 Serre[7]，可能存在整个落在某个非零素数 $p$ 上方的不可约分支。再回顾 Weil 提出的问题：$\operatorname{Hilb}^P_r$ 在各 $s\in\operatorname{Spec}\mathbb Z$ 上的纤维的不可约分支，是否对应于素域的“正则”扩张，即它们是否“相对连通”？这些问题在 Hilbert 概形上，也许比在“Chow 簇”上更容易处理。\footnote{作者后来的\hyperref[orig:8]{C-06页评论}指出，Hartshorne 已研究代数闭域上 Hilbert 概形的连通分支，证明 $\operatorname{Hilb}^P_r$ 连通，并确定哪些 $(r,P)$ 使该概形为空；参见 R. Hartshorne, \textit{Connectedness of the Hilbert scheme}, Harvard 学位论文，1962年。这里保留报告提出问题时的原有叙述。}

\section{若干变体}
\subsection*{a. 拟射影情形}
在(3.1)的条件下，设 $U\subset X$ 为开集，令 $A'$ 为 $A=\underline{\operatorname{Quot}}_{\mathcal F/X/S}$ 的子函子，其中 $A'(S')$ 由 $\mathcal F'$ 的 $S'$-平坦商模 $\mathcal G$ 构成，并要求其支集包含在 $U'$ 中。立即可见，$A'$ 可由表示 $A$ 的预概形 $\operatorname{Quot}_{\mathcal F/X/S}$ 的一个开子集表示。由此可知，定理(3.1)、(3.2)

\origpage{157}{221-19}
也可推广到 $X$ 在 $S$ 上拟射影的情形：把结论中的“射影”同时换成“拟射影”，并把 $\underline{\operatorname{Quot}}_{\mathcal F/X/S}(S')$ 改为 $\mathcal F'$ 的凝聚商层 $\mathcal G$ 的集合，其中 $\mathcal G$ 在 $S'$ 上平坦，且支集在 $S'$ 上固有。

\subsection*{b. 加入其他条件}
一般地，可以对 $\mathcal F'=\mathcal F\otimes_{\mathcal O_S}\mathcal O_{S'}$ 的 $S'$-平坦商层 $\mathcal G$ 加入各种在基变换下稳定的自然条件，从而得到 $\underline{\operatorname{Quot}}_{\mathcal F/X/S}$ 的各种子函子，并寻求表示它们。在许多情形，通常的判别法表明，所得函子仍可由 $\operatorname{Quot}_{\mathcal F/X/S}$ 的开子集表示。特别地，可以加入以下任一种条件：
\begin{enumerate}[label=\arabic*$^\circ$]
\item 对每个 $s'\in S'$，纤维 $X'_{s'}$ 上的诱导模 $\mathcal G_{s'}$ 的所有相伴素闭链的维数，都属于一个给定的整数集合。
\item 当 $\mathcal F=\mathcal O_X$，从而 $\mathcal G$ 对应于 $X'$ 的闭子预概形 $Y$ 时，要求 $Y$ 在 $S'$ 上\term{光滑}{smooth}{simple}（[3]第四章）；或者要求它在 $S'$ 上\term{正规}{normal}{normal}，即纤维 $Y_{s'}$ 在 $k(s')$ 上几何正规，亦即任意基域扩张后仍正规；或者在 $X$ 于 $S$ 上平坦时，要求 $Y$ 是 $X'$ 中关于 $S'$ 的\term{局部完全交}{local complete intersection}{intersection complète locale}，即各 $Y_{s'}$ 在相应的 $X'_{s'}$ 中为局部完全交。
\end{enumerate}
其他条件可以涉及纤维 $X'_{s'}$ 上诱导模 $\mathcal G_{s'}$ 的上同调性质，等等。当然，若若干条件分别由 $\operatorname{Quot}_{\mathcal F/X/S}$ 的开集 $U_i$ 表示，则这些条件的合取由相应开集之交表示。例如，对每个 $S$ 上的 $S'$，考虑 $X'=X\times_SS'$ 中所有为 $S'$ 的给定秩 $r$ 的\term{平展覆盖}{étale covering}{revêtement étale}（[3]第一章）的闭子预概形 $Y$，就得到一个关于 $S'$ 的可表反变函子。

\subsection*{c. 态射与同构的参数概形}
文献[2]第二篇C部分第2节定义的预概形 $\operatorname{Hom}_S(X,Y)$、$\prod_{X/S}(Z/X)$、$\operatorname{Isom}_S(X,Y)$，在适当的射影性假设下存在，并可实现为适当 Hilbert 预概形中的开集。由于
\[
\operatorname{Hom}_S(X,Y)=\prod_{X/S}((X\times_SY)/X),
\]
$\operatorname{Hom}$ 的情形可归结到 $\prod_{X/S}(Z/X)$。注意，对每个 $S$ 上的 $S'$，$Z'=Z\times_SS'$ 在 $X'=X\times_SS'$ 上的截面，与 $Z'$ 中满足如下条件的子预概形 $\Gamma$ 一一对应：由 $Z'\to X'$ 诱导的 $\Gamma\to X'$ 为同构。如果 $Z$ 在 $X$ 上分离，则这样的 $\Gamma$ 必为闭子预概形。由此，

\origpage{158}{221-20}
当 $X$ 在 $S$ 上平坦且固有、$Z$ 在 $S$ 上拟射影时，$\prod_{X/S}(Z/X)$ 存在，且可实现为 $\operatorname{Hilb}_{Z/S}$ 的开子预概形。因此，当 $X$ 在 $S$ 上射影且平坦、$Y$ 在 $S$ 上拟射影时，$\operatorname{Hom}_S(X,Y)$ 存在，且可实现为 $\operatorname{Hilb}_{(X\times_SY)/S}$ 的开子预概形。如果 $X,Y$ 都在 $S$ 上射影，则立即推出 $\operatorname{Isom}_S(X,Y)$ 也存在，并由 $\operatorname{Hom}_S(X,Y)$ 的开集表示。同样，在 $X$ 平坦且射影、$Y$ 拟射影的条件下，表示那些为浸入的 $S'$-态射 $X'\to Y'$ 的子函子的 $S$-预概形 $\operatorname{Imm}_S(X/Y)$，也可由 $\operatorname{Hom}_S(X,Y)$ 的开集表示。

设 $\mathcal L$、$\mathcal M$ 分别为 $X$、$Y$ 上关于 $S$ 甚丰沛的可逆层。由此在 $X\times_SY$ 上得到关于 $S$ 甚丰沛的可逆层 $\mathcal L\otimes_{\mathcal O_S}\mathcal M$。因此，对每个有理系数多项式 $P$，$\operatorname{Hilb}^P_{(X\times_SY)/S}$ 都有定义，且在 $S$ 上拟射影。它在 $\operatorname{Hom}_S(X,Y)$ 中截出一个既开又闭、在 $S$ 上拟射影的部分，记作 $\operatorname{Hom}_S(X,Y)^P$。它在 $S$ 上的截面，就是满足
\[
\chi\!\left((\mathcal L\otimes_{\mathcal O_X}g^*\mathcal M)^{\otimes n}\right)=P(n)
\quad\text{对每个整数 }n
\]
的 $S$-态射 $g:X\to Y$。

这样就推广了 Matsusaka 的定理：一个“极化”射影簇的自同构构成代数群。这里这一断言的含义显然更精确，因为这个群被定义为某个泛问题的解。还应注意，在代数闭域上，以往所考虑的自同构群，是从这里定义的“真正的”群消去幂零元得到的。这解释了为什么旧式构造很难在非完美基域 $k$ 上进行：扩张基域后出现的幂零元理想，不一定“定义在 $k$ 上”。同样的评论也适用于大部分旧式构造。

\origpage{159}{221-21}
\section{Hilbert 概形的微分研究}
这一研究来自下面的结果。

\begin{statement}{命题 5.1}
设 $S$ 为预概形，$S_0$ 为平方为零的拟凝聚理想 $\mathcal I$ 定义的子预概形，$X$ 为 $S$-预概形，$\mathcal F$ 为 $X$ 上的拟凝聚模。记
\[
X_0=X\times_SS_0,\qquad
\mathcal F_0=\mathcal F\otimes_{\mathcal O_S}\mathcal O_{S_0},
\]
并设 $\mathcal G_0=\mathcal F_0/\mathcal H_0$ 为 $\mathcal F_0$ 的拟凝聚商模，在 $S_0$ 上平坦。对每个开集 $U\subset X$，令 $\mathcal E(U)$ 为 $\mathcal F|_U$ 的如下拟凝聚商模 $\mathcal G$ 的集合：$\mathcal G$ 在 $S$ 上平坦，并且 $\mathcal G\otimes_{\mathcal O_S}\mathcal O_{S_0}=\mathcal G_0|_{U\cap X_0}$。随着 $U$ 变化，这些集合构成 $X$ 上的层 $\mathcal E$ 的截面。则群层
\[
\mathcal A=\mathcal{H}om_{\mathcal O_{X_0}}
\left(\mathcal H_0,\mathcal G_0\otimes_{\mathcal O_{S_0}}\mathcal I\right)
\]
自然作用于 $\mathcal E$，使后者成为一个“形式上在 $\mathcal A$ 下主齐性的层”：亦即，对每个开集 $U\subset X$，$\mathcal E(U)$ 或者为空，或者是 $\mathcal A(U)$ 下的\term{主齐性集}{principal homogeneous set}{ensemble principal homogène}。
\end{statement}
由此得到：
\begin{statement}{推论 5.2}
假设在 $X$ 上局部地，存在把 $\mathcal G_0$ 延拓为 $\mathcal F$ 的一个 $S$-平坦商层 $\mathcal G$ 的方法，即层 $\mathcal E$ 的各茎均非空。则存在自然的\term{障碍类}{obstruction class}{classe d'obstruction}
\[
c(\mathcal G_0)\in H^1(X,\mathcal A),
\]
其消失是存在上述全局延拓的充分必要条件。若该类为零，则所有可能延拓组成的集合 $\mathcal E(X)$ 是
\[
\mathcal A(X)=\operatorname{Hom}_{\mathcal O_{X_0}}
\left(\mathcal H_0,\mathcal G_0\otimes_{\mathcal O_{S_0}}\mathcal I\right)
\]
下的主齐性集。
\end{statement}
特别地，如果 $H^1(X,\mathcal A)=0$，全局延拓必然存在。

\begin{statement}{推论 5.3}
假设 $Q=\operatorname{Quot}_{\mathcal F/X/S}$ 存在（参见第4节a）；例如，$S$ 局部诺特，$X$ 在 $S$ 上拟射影，且 $\mathcal F$ 凝聚。设 $x\in Q$，$K=k(x)$ 为某个 $k(s)$（$s\in S$）的扩张。

\origpage{160}{221-22}
于是 $x$ 由 $K$-预概形 $X_K$ 上的凝聚商模 $\mathcal G_0=\mathcal F_0/\mathcal H_0$ 定义，其中 $\mathcal F_0=\mathcal F_K$。在 $X_K$ 上令
\[
\mathcal A=\mathcal{H}om_{\mathcal O_{X_K}}(\mathcal H_0,\mathcal G_0).
\]
则纤维在 $x$ 处的\term{Zariski 切空间}{Zariski tangent space}{espace tangent de Zariski}，自然同构于 $H^0(X_K,\mathcal A)$。这里切空间是在基扩张 $Q_K$ 的相应 $K$-点处计算的，即 $\mathfrak m/\mathfrak m^2$ 的 $K$-对偶，其中 $\mathfrak m$ 为 $\mathcal O_{Q_K,x}$ 的极大理想。
\end{statement}
这个 Zariski 切空间公式可以推广。对给定的 $S$-态射 $g:S'\to Q$，或者说 $Q'=Q\times_SS'$ 在 $S'$ 上的截面 $g'$，它给出模
\[
\Omega=g^*\Omega^1_{Q/S}=g'^*(\mathcal J/\mathcal J^2)
\]
的一种刻画，其中 $\mathcal J$ 为 $Q'$ 上定义截面 $g'$ 的理想。具体地，对 $S'$ 上凝聚模 $\mathcal M$，有关于 $\mathcal M$ 函子性的公式
\[
\operatorname{Hom}_{\mathcal O_{S'}}(\Omega,\mathcal M)
\simeq H^0(X',\mathcal A),
\]
这里 $X'=X\times_SS'$，且
\[
\mathcal A=\mathcal{H}om_{\mathcal O_{X'}}
(\mathcal H,\mathcal G\otimes_{\mathcal O_{S'}}\mathcal M),
\]
其中 $\mathcal G=\mathcal F'/\mathcal H$ 为对应于 $g$ 的 $\mathcal F'=\mathcal F\otimes_{\mathcal O_S}\mathcal O_{S'}$ 的商模。只需在(5.1)中，以 $S'$ 替换 $S_0$，并以预概形
\[
D(\mathcal M)=(S',\mathcal O_{S'}\oplus\mathcal M)
\]
替换 $S$ 即可；这里把 $\mathcal M$ 视为平方为零的理想。

当(5.1)中 $\mathcal F=\mathcal O_X$ 时，给定 $\mathcal G_0$ 就相当于给定 $X_0$ 的一个 $S_0$-平坦闭子预概形 $Y_0$，其定义理想为 $\mathcal J_0=\mathcal H_0$。此时前面的公式给出
\[
\mathcal A=\mathcal{H}om_{\mathcal O_{X_0}}
\left(\mathcal J_0/\mathcal J_0^2,
\mathcal O_{Y_0}\otimes_{\mathcal O_{S_0}}\mathcal I\right).
\]
其中 $\mathcal J_0/\mathcal J_0^2$ 是 $Y_0$ 在 $X_0$ 中的\term{余法层}{conormal sheaf}{faisceau conormal}，也记作

\origpage{161}{221-23}
$\mathcal N_{Y_0/X_0}$。这时适合把 $\mathcal A$ 看作 $Y_0$ 上的模，并在 $Y_0$ 上计算 $H^0$ 与 $H^1$。如果 $Y_0$ 在 $X_0$ 中局部为完全交，且 $X$ 在 $S$ 上平坦，则(5.1)中的局部延拓必然存在。另一方面，$\mathcal J_0/\mathcal J_0^2$ 在 $Y_0$ 上局部自由，可以写成
\[
\mathcal A=\mathcal N^\vee_{Y_0/X_0}
\otimes_{\mathcal O_{S_0}}\mathcal I.
\]
右侧的第一个因子是 $Y_0$ 在 $X_0$ 中的\term{法层}{normal sheaf}{faisceau normal}。\translatornote{原报告这一页把对偶法层的下标次序写作 $X_0/Y_0$；这里统一为与余法层相同的 $Y_0/X_0$，上标 $\vee$ 表示取对偶。}利用基本的光滑性判别法（[3]第三章3.1），例如得到：

\begin{statement}{推论 5.4}
在(5.3)的条件下，再假设 $\mathcal F=\mathcal O_X$、$X$ 在 $S$ 上平坦，且对应于 $\mathcal G_0$ 的闭子预概形 $Y_0\subset X_K$ 局部为完全交。那么，$Q_K$ 在 $x$ 处的 Zariski 切空间自然同构于 $H^0(Y_0,\mathcal N^\vee_{Y_0/X_K})$。若
\[
H^1(Y_0,\mathcal N^\vee_{Y_0/X_K})=0,
\]
则 Hilbert 预概形 $Q$ 在点 $x$ 处关于 $S$ 光滑。这里 $\mathcal N^\vee_{Y_0/X_K}$ 为 $Y_0$ 在 $X_K$ 中的法层。
\end{statement}

\subsection*{注 5.5}
这个结论特别适用于 $Y_0$ 是 $X_K$ 中由一个方程定义的完全交，即为\term{有效 Cartier 除子}{effective Cartier divisor}{diviseur de Cartier positif}的情形。此时，$\mathcal N^\vee_{Y_0/X_K}$ 同构于 $X_K$ 上由除子 $Y_0$ 定义的可逆层 $\mathcal J_0^{-1}$ 在 $Y_0$ 上诱导的层。研究非奇异射影簇 $X$ 上的有效除子族时，尤其会遇到这种情形。$Q$（或者其中对应于除子的开集 $D$）在点 $x$ 处的 Zariski 切空间与 $H^0(Y_0,\mathcal N^\vee_{Y_0/X_K})$ 之间的同构，在经典代数几何中以从前者到后者的“\term{特征同态}{characteristic homomorphism}{homomorphisme caractéristique}”著称。

在旧理论中，这个同态仅当 $x$ 是除子的某个“完备连续族”的参数簇 $T$ 的光滑点时才有定义。从我们的角度看，$T$ 就是赋予诱导约化结构的概形 $D$ 的一个不可约分支。因此，$T$ 在 $x$ 处的切空间是 $D$ 在 $x$ 处切空间的子空间。旧理论的特征同态确实是单射，但只有加入其他条件才是满射，例如 $D$ 在 $x$ 处整时。事实上，Zappa[8]构造了一个例子，

\origpage{162}{221-24}
其中 $X$ 是复数域上的非奇异射影曲面，但即使在 $T$ 的泛点处，特征同态也不满射。也就是说，即使在所考察不可约分支的泛点，$D$ 也不是整的。这特别鲜明地说明：即使为了理解最经典的曲面理论中的现象，也必须引入带幂零元的簇。\footnote{作者在\hyperref[orig:8]{C-06页评论}中进一步指出，Mumford 已构造 $\mathbb P^3_{\mathbb C}$ 的 Hilbert 概形的一个不可约分支，其泛点非约化，一般点参数化次数14、亏格24的非奇异曲线。沿其中一条曲线作爆破，还得到一个正则射影三维概形，其形式模概形在泛点非约化；等价地，其复解析意义下的局部模簇在各点均非约化。}

\subsection*{5.6}
(5.4)给出了一个光滑性判别法，尤其适用于除子概形。Kodaira 在[6]中给出了另一判别法，即
\[
H^1(X_K,\mathcal L)=0,
\]
其中 $\mathcal L=\mathcal J_0^{-1}$ 是 $X_K$ 上由除子 $Y_0$ 定义的可逆层。这个判别法适用于 $S$ 为特征零域的谱的情形；[6]在基域为 $\mathbb C$ 时以超越方法证明了它。

这里指出，一般地，即使 $S$ 任意，Kodaira 条件仍足以保证从除子预概形到 $X/S$ 的 Picard 预概形的自然态射
\[
D\longrightarrow\operatorname{Pic}_{X/S}
\]
在所考察的点 $x$ 处光滑。一旦知道 $\operatorname{Pic}_{X/S}$ 存在，用通常的光滑性判别法即可验证这一点。因此，如果 $\operatorname{Pic}_{X/S}$ 在 $x$ 的像点处关于 $S$ 也光滑，例如整个 $\operatorname{Pic}_{X/S}$ 在 $S$ 上光滑，那么 $D$ 在 $x$ 处关于 $S$ 光滑。另一方面，Cartier 证明了：特征零域 $k$ 上每个局部有限型群预概形都在 $k$ 上光滑。结合这两个结果，就重新得到 Kodaira 的结论。

这些观察还说明，在特征 $p>0$ 的域 $K$ 上，如果 $\operatorname{Pic}_{X/K}$ 不光滑（例如 $X$ 为 Igusa 曲面时），条件 $H^1(X_K,\mathcal L)=0$ 反而意味着 $D$ 在 $x$ 处不光滑；若 $K$ 代数闭，它甚至在 $x$ 处非约化。

最后再指出一个结果，它在纤维空间的微分研究中起着重要作用。
\begin{statement}{命题 5.7}
设 $X$ 在局部诺特预概形 $S$ 上有限且平坦，$Z$ 为 $X$ 上的预概形，并假设 $\prod_{X/S}(Z/X)$ 存在；当 $Z$ 在 $X$ 上拟射影时，这一存在性成立。若 $Z$ 在 $X$ 上光滑，则 $\prod_{X/S}(Z/X)$ 在 $S$ 上光滑。
\end{statement}
这直接来自定义及通常的光滑性判别法（[3]第三章3.1）。注意，当 $X$ 在 $S$ 上有限且平坦时，$\prod_{X/S}(Z/X)$ 的存在性问题可以用十分初等的方法处理，

\origpage{163}{221-25}
无需使用 Hilbert 概形理论。例如，若 $X$ 在 $S$ 上\term{根式}{radicial}{radiciel}，那么无需对 $Z$ 加任何限制，$\prod_{X/S}(Z/X)$ 就存在。

例如，设 $T$ 为 $S$-预概形，$T_n$ 为对角线在 $T\times_ST$ 中的“$n$ 阶\term{无穷小邻域}{infinitesimal neighbourhood}{voisinage infinitésimal}”，并带有两个投影诱导的态射 $p_1,p_2:T_n\to T$。利用 $p_1$，把 $T_n$ 看作 $T$ 上的有限预概形，再假设它在 $T$ 上平坦；当 $T$ 在 $S$ 上光滑时，这一条件成立。对每个 $T$ 上的预概形 $X$，令
\[
(X/T/S)^{(n)}=\prod_{T_n/T}\bigl(p_2^*(X/T)/T_n\bigr).
\]
这是一个 $T$ 上的预概形，称为 $X$ 在 $T$ 上、相对于 $S$ 的\term{$n$ 阶截面芽丛}{bundle of germs of sections of order n}{fibré des germes de sections d'ordre n}。它函子性地依赖于 $X$；若 $X$ 在 $T$ 上光滑，它也在 $T$ 上光滑。\translatornote{原式把限制标量的下标分母印为 $S$。按前一句给定的有限平坦态射 $p_1:T_n\to T$，此处应为 $T$，否则所得对象不是文中所说的 $T$ 上截面芽丛。例如 $S=\operatorname{Spec}k$、$T=\mathbb A^1_k$、$X=\mathbb A^1_T$ 时，正确构造的纤维记录 $n+1$ 个 Taylor 系数，得到 $T$ 上的仿射 $(n+1)$ 维空间。}

\section{与范数概念及对称积的关系}
设 $S$ 为预概形，$X,Y$ 为 $S$-预概形，
\[
u:(X/S)^n\longrightarrow Y
\]
为从 $X/S$ 的 $n$ 次笛卡尔幂到 $Y$ 的对称 $S$-态射。为简便起见，假设 $S$ 局部诺特，$X,Y$ 在 $S$ 上有限型。给定 $X$ 上支集在 $S$ 上有限、在 $S$ 上平坦且相对秩为 $n$ 的凝聚模 $\mathcal F$，即 $f_*(\mathcal F)$ 是 $S$ 上局部自由秩 $n$ 的模，可以自然地对应一个 $Y$ 在 $S$ 上的截面：
\[
\mathcal N^u_{X/S}(\mathcal F)\in\Gamma(Y/S).
\]
这里不详细给出定义，只指出所得形式体系是通常\term{范数}{norm}{norme}与\term{迹}{trace}{trace}的形式体系的自然推广。

如果 $X$ 在 $S$ 上的 $n$ 次\term{对称幂}{symmetric power}{puissance symétrique}存在，例如对称群 $\mathfrak S_n$ 作用于 $(X/S)^n$ 的每条轨道都包含在某个仿射开集中，就可以令 $Y$ 为该对称幂 $\operatorname{Symm}^n_S(X)$，从而得到自然元素
\[
\mathcal N_{X/S}(\mathcal F)\in\Gamma(\operatorname{Symm}^n_S(X)/S),
\]

\origpage{164}{221-26}
它可以恢复前面的所有 $\mathcal N^u_{X/S}(\mathcal F)$。另一个重要情形是 $X$ 为 $S$ 上的交换幺半群，$Y=X$，态射 $u$ 由 $X$ 的合成律给出。此时，把与 $X$ 上模 $\mathcal F$ 对应的 $X$ 在 $S$ 上的截面，简记为 $\mathcal N(\mathcal F)$。

现设 $X$ 上有凝聚模 $\mathcal F$，使 $\operatorname{Quot}_{\mathcal F/X/S}$ 存在；或者至少如下函子 $\underline{\operatorname{Quot}}^n_{\mathcal F/X/S}$ 可由 $S$-预概形 $\operatorname{Quot}^n_{\mathcal F/X/S}$ 表示：它把 $S$ 上的每个 $S'$ 对应到 $\mathcal F'=\mathcal F\otimes_{\mathcal O_S}\mathcal O_{S'}$ 的凝聚商层 $\mathcal M$ 的集合，要求 $\mathcal M$ 在 $S'$ 上平坦，支集有限，相对秩为 $n$。当 $X$ 在 $S$ 上拟射影时，$\operatorname{Quot}^n_{\mathcal F/X/S}$ 确实存在；用第3节的记号，它就是 $P$ 为常数多项式 $n$ 时的 $\operatorname{Quot}^P_{\mathcal F/X/S}$。由于 $\mathcal N^u_{X/S}(\mathcal M)$ 的形成与基变换相容，得到自然态射
\[
\mathcal N^u_{X/S}:\operatorname{Quot}^n_{\mathcal F/X/S}\longrightarrow Y.
\]
特别地，若 $X$ 在 $S$ 上的 $n$ 次对称幂存在，则有
\[
\mathcal N_{X/S}:\operatorname{Quot}^n_{\mathcal F/X/S}
\longrightarrow\operatorname{Symm}^n_S(X).
\]
最重要的是 $\mathcal F=\mathcal O_X$ 的情形，它给出态射
\[
\mathcal N_{X/S}:\operatorname{Hilb}^n_{X/S}
\longrightarrow\operatorname{Symm}^n_S(X).
\]
当 $n=0$ 或 $n=1$ 时，这显然为同构。但当 $n>1$ 时，即使 $S$ 是域 $k$ 的谱，$X$ 在 $S$ 上光滑，一般也既不是同构，甚至不一定是单射。因为知道 $X$ 的一个零维子概形所定义的闭链，并不足以确定该子概形。例如，给定闭点 $x\in X$，在局部环 $\mathcal O_{X,x}$ 中取一个关于极大理想的准素理想 $\mathfrak a$；知道对应闭链，在这里也就是知道 $\mathfrak a$ 在 $\mathcal O_{X,x}$ 中的 $k$-余维，并不能确定 $\mathfrak a$。对一般的 $S$，只能作如下断言。

\noindent\textbf{a.} 若 $X$ 在 $S$ 上光滑，则范数态射在两个开集之间诱导同构：一边是 $\operatorname{Hilb}^n_{X/S}$ 中分类 $X$ 内秩 $n$ 平展覆盖的开集（参见第4节b）；另一边是 $\operatorname{Symm}^n_S(X)$ 中对应于没有重分支的 $n$ 次闭链的开集。

\origpage{165}{221-27}
\noindent\textbf{b.} 若进一步假设 $X$ 在 $S$ 上的相对维数为1，则范数态射甚至给出整个 $\operatorname{Hilb}^n_{X/S}$ 与 $\operatorname{Symm}^n_S(X)$ 之间的同构。

后一个事实源于：非奇异代数曲线上的零维子概形，可由其对应闭链确定。更一般地，非奇异代数概形上的有效 Cartier 除子也有相同性质。在这个极为特殊的情形，也许 Chow 簇与 Hilbert 概形确实给出同一个对象。

\section{补充与问题}
正如 J.-P. Serre 所指出的，由 Nagata 的一个著名例子可知，存在如下对象：$S=\operatorname{Spec}k$；$S$-概形 $S'=\operatorname{Spec}k'$，其中 $k'/k$ 为二次扩张；以及 $S'$ 上一个三维固有光滑但非射影的概形 $X$，使得 $\prod_{S'/S}(X/S')$ 不存在。\translatornote{原文此处的括号内对象印作 $X/S$；其下标及上下文表明这里是沿 $S'\to S$ 对 $S'$-概形 $X$ 作限制标量，因此写作 $X/S'$。}这更意味着 Hilbert 概形 $\operatorname{Hilb}_{X/S}$ 不存在；甚至表示 $X$ 中秩2的 $S$ 上平展覆盖的 $k$-概形也不存在，更不用说 $X$ 的二次对称幂了，参见上一节。因此，代数几何中的非射影构造确实受到相当严重的限制。不过，解析几何中可能不会出现这些限制，就像形式几何中没有这些限制一样（参见[2]第二篇）。

相反，若 $X$ 是域 $k$ 的谱 $S$ 上的固有概形，$Z$ 在 $X$ 上拟射影，则 $\prod_{X/S}(Z/X)$ 存在，且是一个概形，由一列在 $S$ 上拟射影的概形取不交并而成，如同射影情形(3.1)。证明时，取一个在 $S$ 上射影、支配 $X$ 的概形 $X'$，就能归结到 $X$ 本身射影的情形。这里不详述证明；它还使用了[2]第一篇A部分第2节(b)指出的有限态射分解结果。这个方法成功的关键是，$S$ 为域的谱，因而 Chow 引理中出现的 $X'$ 自动在 $S$ 上平坦。

如果不对 $S$ 作额外假设，只要求 $X$ 在 $S$ 上固有且平坦、$Z$ 在 $X$ 上拟射影，我不知道上述结果是否仍成立。在应用中一个重要情形是 $Z$ 为 $X$ 的闭子概形；如果此时 $\prod_{X/S}(Z/X)$ 存在，它必为 $S$ 的闭子概形。当 $X$ 在 $S$ 上射影时，可以

\origpage{166}{221-28}
很简单地直接构造它，无需借助 Hilbert 概形理论。这一方法更一般地还说明，若 $Z$ 在 $X$ 上仿射，则 $\prod_{X/S}(Z/X)$ 存在，并在 $S$ 上仿射。该方法还表明：如果 $X$ 在 $S$ 上固有且平坦，不必射影，那么对 $X$ 上每个局部平凡向量丛 $Z$，$\prod_{X/S}(Z/X)$ 都存在，并为 $S$ 上的线性丛。\translatornote{此处的 \textit{fibré vectoriel} 应按允许纤维秩变化的线性概形理解，不能无条件理解为现代意义下局部平凡的向量丛；后者需要相应直接像的局部自由性及基变换条件。例如取椭圆曲线 $E$，令 $S=\operatorname{Pic}^0(E)$、$X=E\times S$，$\mathcal L$ 为归一化 Poincaré 线丛，$Z$ 为其全空间。其截面参数对象在 $s=0$ 处的纤维为 $H^0(E,\mathcal O_E)\simeq k$，在非平凡零次线丛对应的点处为0；因此它不是局部恒秩的向量丛。}希望能够重新整理并统一这些结果。

\section*{参考文献}
\begin{description}
\item[{[1]}] A. Borel、J.-P. Serre：《Riemann--Roch 定理》（\textit{Le théorème de Riemann--Roch}），\textit{Bull. Soc. math. France}，第86卷，1958年，第97--136页。
\item[{[2]}] Alexander Grothendieck：《代数几何中的构造技术与存在定理》第一、二、三篇（\textit{Technique de construction et théorèmes d'existence en géométrie algébrique, I, II, III}），\textit{Séminaire Bourbaki}，第13卷，1959/60年，第190号（29页）、第195号（22页）；第14卷，1960/61年，第212号（20页）。
\item[{[3]}] Alexander Grothendieck：《代数几何讨论班》第一至四部分（\textit{Séminaire de Géométrie algébrique, I, II, III, IV}），巴黎，高等科学研究所，1960/61年，油印本。
\item[{[4]}] Alexander Grothendieck：《解析几何中的构造技术》第四、五篇（\textit{Technique de construction en géométrie analytique, IV, V}），\textit{Séminaire Cartan}，第13卷，1960/61年，第11、12号。
\item[{[5]}] A. Grothendieck、J. Dieudonné：《代数几何原理》（\textit{Éléments de géométrie algébrique}），第一章《概形的语言》（\textit{Le langage des schémas}），巴黎，法国大学出版社，1960年（高等科学研究所《数学出版物》第4号）；第二、三、四章待刊。
\item[{[6]}] K. Kodaira：《完备连续系的特征线性系》（\textit{Characteristic linear systems of complete continuous systems}），\textit{Amer. J. of Math.}，第78卷，1956年，第716--744页。
\item[{[7]}] Jean-Pierre Serre：《特征 $p$ 中不能提升到特征零的射影簇的例子》（\textit{Exemples de variétés projectives en caractéristique p non relevables en caractéristique 0}），\textit{Proc. Nat. Acad. Sc. U.S.A.}，第47卷，1961年，第108--109页。
\item[{[8]}] G. Zappa：《代数曲面上存在其曲线的特征级数不完备的无限完备连续系》（\textit{Sull'esistenza, sopra la superficie algebriche, di sistemi continui completi infiniti, la cui curva è a serie caratteristica incompleta}），\textit{Pont. Acad. Sc. Acta}，第9卷，1945年，第91--93页。
\end{description}

\section*{补遗}
\begin{center}（校样校订时补入）\end{center}
现在已经知道，[2]第三篇第8节的猜想是错误的：即使对特征零域上的非奇异簇，关于商的存在性及拟射影性的猜想都不成立；甚至群 $G$ 的作用具有闭图形时，结论仍然如此。
